\chapter{分数：比较与加减}\label{ch:g7:fractions}

六年级已将\omterm{def:g6:fractions:def}{分数}作为份数与精确的商引入
（见\cref{ch:g6:fractions}）。本章学习对\omterm{def:g6:fractions:def}{分数}进行\emph{计算}，
包括识别相等的\omterm{def:g6:fractions:def}{分数}、比较大小、加减运算——
目前只讨论\omterm{def:g4:fractions:def}{分母}较为简单的情形，
一般情形见\cref{ch:g9:fractions}。

\section{相等的分数}

\begin{theorem}[相等的分数]\label{thm:g7:fractions:equal}
将\omterm{def:g6:fractions:def}{分数}的\omterm{def:g4:fractions:def}{分子}和\omterm{def:g4:fractions:def}{分母}同乘或同除以同一个非零数，
\omterm{def:g6:fractions:def}{分数}的值保持不变：
\[
\frac{a}{b} = \frac{a \times k}{b \times k} .
\]
\end{theorem}
\admitted

\begin{example}\label{ex:g7:fractions:equal}
$\dfrac{3}{4} = \dfrac{3 \times 5}{4 \times 5} = \dfrac{15}{20}$，反过来则有
即 $\dfrac{42}{30} = \dfrac{42 \div 6}{30 \div 6}
= \dfrac{7}{5}$（用 $6$ 约分）。

要判断 $\dfrac{8}{12}$ 与 $\dfrac{10}{15}$ 是否相等，
只需将两者化简：$\dfrac{8}{12} = \dfrac{2}{3}$ 且
$\dfrac{10}{15} = \dfrac{2}{3}$——两者相等。
\end{example}

\section{比较分数的大小}

\begin{method}[比较两个分数的大小]\label{met:g7:fractions:compare}
\begin{enumerate}
  \item \emph{\omterm{def:g4:fractions:def}{分母}相同}时，比较\omterm{def:g4:fractions:def}{分子}：
        $\frac{5}{9} < \frac{7}{9}$。
  \item \emph{一个\omterm{def:g4:fractions:def}{分母}是另一个的\omterm{def:g5:division:multiple}{倍数}}时，将
        \omterm{def:g4:fractions:def}{分母}较小的\omterm{def:g6:fractions:def}{分数}改写后再比较：
        $\frac{5}{6}$ 与 $\frac{7}{12}$ 比较：$\frac56 = \frac{10}{12} >
        \frac{7}{12}$。
  \item 若可能，\emph{与 $1$ 或 $\frac12$ 比较}：
        $\frac{9}{8} > 1 > \frac{7}{9}$，由此立即判定
        $\frac98$ 与 $\frac79$ 的大小。
\end{enumerate}
\end{method}

\begin{omfigure}
\begin{tikzpicture}[scale=0.62]
  \foreach \i in {0,...,5} \draw[omExo] (\i*2,0) rectangle (\i*2+2,1.1);
  \foreach \i in {0,...,4} \fill[omDef!35] (\i*2+0.05,0.05)
    rectangle (\i*2+1.95,1.05);
  \draw[thick] (0,0) rectangle (12,1.1);
  \node[right, font=\small] at (12.3,0.55) {$\frac56$};
  \begin{scope}[yshift=-1.9cm]
  \foreach \i in {0,...,11} \draw[omExo] (\i,0) rectangle (\i+1,1.1);
  \foreach \i in {0,...,6} \fill[omThm!35] (\i+0.05,0.05)
    rectangle (\i+0.95,1.05);
  \draw[thick] (0,0) rectangle (12,1.1);
  \node[right, font=\small] at (12.3,0.55) {$\frac{7}{12}$};
  \end{scope}
\end{tikzpicture}

{\small 用双条形图比较 $\frac56$ 与 $\frac{7}{12}$：
将六分之一各切成两半，可得 $\frac56 = \frac{10}{12}$，多于
$\frac{7}{12}$。}
\end{omfigure}

\section{分数的加减}

\begin{theorem}[同分母分数的加减]\label{thm:g7:fractions:add}
\omterm{def:g4:fractions:def}{分母}相同的\omterm{def:g6:fractions:def}{分数}相加（或相减）时，
\omterm{def:g4:fractions:def}{分子}相加（或相减）即可：
\[
\frac{a}{d} + \frac{b}{d} = \frac{a + b}{d},
\qquad
\frac{a}{d} - \frac{b}{d} = \frac{a - b}{d} .
\]
\end{theorem}

\begin{proof}[理由]
$a$ 份大小为 $\frac1d$ 的部分与 $b$ 份同样大小的部分合起来，
共得 $a + b$ 份这样大小的部分。
\end{proof}

\begin{method}[分母不同的分数加减]\label{met:g7:fractions:add}
当一个\omterm{def:g4:fractions:def}{分母}是另一个的\omterm{def:g5:division:multiple}{倍数}时：
\begin{enumerate}
  \item 将\omterm{def:g4:fractions:def}{分母}较小的\omterm{def:g6:fractions:def}{分数}改写，使两者\omterm{def:g4:fractions:def}{分母}一致，即
        都为较大的那个\omterm{def:g4:fractions:def}{分母}（见\cref{thm:g7:fractions:equal}）；
  \item \omterm{def:g4:fractions:def}{分子}相加或相减；
  \item 若可能，化简结果。
\end{enumerate}
\end{method}

\begin{example}\label{ex:g7:fractions:add}
计算 $\dfrac{3}{4} + \dfrac{5}{8}$，分步进行：
\begin{enumerate}
  \item $8$ 是 $4$ 的\omterm{def:g5:division:multiple}{倍数}，将
        $\dfrac34 = \dfrac{3 \times 2}{4 \times 2} = \dfrac{6}{8}$；
  \item 相加：$\dfrac{6}{8} + \dfrac{5}{8} = \dfrac{11}{8}$；
  \item 无需化简，结果为 $\dfrac{11}{8}$，即\
        $1 + \dfrac38$。
\end{enumerate}
另一例：$2 - \dfrac{4}{5} = \dfrac{10}{5} - \dfrac{4}{5} =
\dfrac{6}{5}$（先把整数写成以 $5$ 为\omterm{def:g4:fractions:def}{分母}的\omterm{def:g6:fractions:def}{分数}
）。
\end{example}

\begin{example}[经典陷阱]\label{ex:g7:fractions:trap}
$\dfrac12 + \dfrac13$ 并不等于\emph{不}$\dfrac{2}{5}$！
分别把\omterm{def:g4:fractions:def}{分子}与\omterm{def:g4:fractions:def}{分母}相加是错误的——答案 $\frac25$ 会
小于 $\frac12$，而\emph{加上}一个正数后，
结果应比 $\frac12$ 更大，显然荒谬。（正确的和 $\frac56$ 需要
以 $6$ 为公\omterm{def:g4:fractions:def}{分母}；一般方法见\cref{ch:g9:fractions}。）
\end{example}

\section{分数乘以一个数}

\begin{proposition}[分数乘以一个数]\label{prop:g7:fractions:mult}
设 $k$ 为一个数：
\[
k \times \frac{a}{b} = \frac{k \times a}{b} .
\]
求某个量\emph{的} $\frac ab$，即是乘以
$\frac ab$。
\end{proposition}
\admitted

\begin{example}\label{ex:g7:fractions:mult}
$6 \times \dfrac{2}{3} = \dfrac{12}{3} = 4$：六的三分之二是
四个整体。“$\frac34$ 的 $60$ 欧元”即
$\frac34 \times 60 = \frac{180}{4} = 45$ 欧元，与先除以 $4$ 再
乘以 $3$ 所得的结果相同
（见\cref{met:g6:fractions:ofquantity}）。
\end{example}

\section{练习}

\begin{exercise}[$\star$]\label{exo:g7:fractions:1}
填空：$\dfrac{2}{5} = \dfrac{?}{20}$；\quad
$\dfrac{18}{24} = \dfrac{3}{?}$；\quad
$\dfrac{7}{3} = \dfrac{28}{?}$；\quad
$\dfrac{?}{9} = \dfrac{20}{36}$。
\end{exercise}

\begin{exercise}[$\star$]\label{exo:g7:fractions:2}
尽可能约分：
$\dfrac{12}{16}$；\quad $\dfrac{30}{45}$；\quad $\dfrac{27}{36}$；
\quad $\dfrac{48}{12}$。
\end{exercise}

\begin{exercise}[$\star$]\label{exo:g7:fractions:3}
$\dfrac{15}{35}$ 与 $\dfrac{21}{49}$ 相等吗？$\dfrac{16}{28}$
与 $\dfrac{20}{36}$ 呢？通过约分说明理由。
\end{exercise}

\begin{exercise}[$\star$]\label{exo:g7:fractions:4}
比较大小（写出推理过程，而不仅写出答案）：
\[
\frac{7}{11} \text{ 与 } \frac{9}{11}; \qquad
\frac{5}{6} \text{ 与 } \frac{11}{18}; \qquad
\frac{13}{12} \text{ 与 } \frac{19}{20}.
\]
\end{exercise}

\begin{exercise}[$\star$]\label{exo:g7:fractions:5}
计算，并在可能时约分：
\[
\frac{5}{9} + \frac{2}{9}, \qquad
\frac{11}{7} - \frac{4}{7}, \qquad
\frac{5}{8} + \frac{7}{8} .
\]
\end{exercise}

\begin{exercise}[$\star$]\label{exo:g7:fractions:6}
利用\cref{met:g7:fractions:add}计算：
\[
\frac{2}{3} + \frac{5}{6}, \qquad
\frac{7}{10} - \frac{2}{5}, \qquad
\frac{3}{4} + \frac{5}{12} .
\]
\end{exercise}

\begin{exercise}[$\star$]\label{exo:g7:fractions:7}
计算：
\[
3 - \frac{5}{4}, \qquad
1 + \frac{3}{8}, \qquad
2 - \frac{7}{6} .
\]
\end{exercise}

\begin{exercise}[$\star$]\label{exo:g7:fractions:8}
计算：
\[
5 \times \frac{3}{10}, \qquad
\frac{7}{8} \times 4, \qquad
\frac{2}{3} \text{ 的 } 45 .
\]
\end{exercise}

\begin{exercise}[$\star\star$]\label{exo:g7:fractions:9}
马克吃了一个披萨的 $\frac14$，朱莉吃了同一个披萨的 $\frac38$。
他们一共吃了披萨的几分之几？
还剩下几分之几？
\end{exercise}

\begin{exercise}[$\star\star$]\label{exo:g7:fractions:10}
一瓶果汁有 $\frac34$ 升。莉亚
两次各倒出 $\frac16$ 升。还剩多少果汁？（公\omterm{def:g4:fractions:def}{分母}为 $12$。）
\end{exercise}

\begin{exercise}[$\star\star$]\label{exo:g7:fractions:11}
利用\cref{ex:g7:fractions:trap}中的论证，说明为何
$\frac35 + \frac12$ 不可能等于 $\frac{4}{7}$，无需计算
正确的和。
\end{exercise}

\begin{exercise}[$\star\star\star$]\label{exo:g7:fractions:12}
计算下面的和
\[
\frac12 + \frac14 + \frac18 + \frac{1}{16}
\]
按步骤（从左到右）进行。观察部分和的规律：
随着越来越多的项
$\frac{1}{32}, \frac{1}{64}, \dots$ 被加入，它们趋近于什么数？
\end{exercise}

\section{问题：藏在一小节乐曲里的分数}

\begin{problem}[{周末问题——音符时值是必须恰好相加的分数，
节奏计数隐藏着一个著名数列，
奇偶性主宰巴尔干节拍}]
\label{pb:g7:fractions:1}
乐谱是用\omterm{def:g6:fractions:def}{分数}记写的。一个\emph{全音符}持续 $1$
小节（在常见的“四四拍”中）；一个\emph{二分音符}持续 $\frac12$；
随后是\emph{四分音符} $\frac14$、
\emph{八分音符} $\frac18$ 和\emph{十六分音符}
$\frac{1}{16}$。一条不可违反的规则：\emph{一小节内的各时值
之和必须恰好等于该小节的总时值}——在四四拍中即 $1$。
下面的每个问题都是本章的算术
（\cref{thm:g7:fractions:add}，\cref{met:g7:fractions:add}）的乐谱版本；
沿途你还将遇到一个计数数列，
它由印度学者在一千年前发现。

\medskip
\noindent\textbf{第一部分——时值必须凑齐的音符。}
\begin{enumerate}
  \item 检验下列每个小节在四四拍中是否合法：
        (a) 二分音符 $+$ 四分音符 $+$ 四分音符；(b) 四分音符 $+$
        八分音符 $+$ 八分音符 $+$ 二分音符。
  \item 一位抄谱员写下了：二分音符 $+$ 四分音符 $+$ 八分音符。
        写下的部分占该小节的几分之几？
        还需要哪一个单独的音符才能将其补全？
  \item 多少个十六分音符能填满一整个小节？
        多少个十六分音符组成一个八分音符？
  \item 音符后面的\emph{附点}使该音符的时值增加其\omterm{def:g3:division:half}{一半}：附点二分音符是 $\frac12 + \frac14$，
        附点四分音符是 $\frac14 + \frac18$，
        请计算这两个时值。
  \item 华尔兹以三四拍记写，每小节的总时值必须
        为 $\frac34$。检验一个附点二分音符能否恰好填满一个华尔兹小节，
        并另写出两个合法的华尔兹小节，
        仅使用四分音符和八分音符。
\end{enumerate}

\medskip
\noindent\textbf{第二部分——抄谱员的工坊。}
\begin{enumerate}[resume]
  \item 一个四四拍小节中包含：附点四分音符 $+$ 八分音符
        $+$ 四分音符。还缺什么？
  \item 附点八分音符与四分音符，哪一个时值更长？
        请借助十六分音符进行比较。
  \item \emph{连音线}把若干时值连成一个较长的声音：
        一个二分音符与一个八分音符相连，时值有多长？
        这相当于多少个十六分音符？
  \item 一个小节中各音符的时值之和为 $\frac{11}{16}$，余下的部分
        是无声的（\emph{休止符}），
        这一休止符的时值是多少？
  \item 一位鼓手要填满一个四分音符（占小节的 $\frac{4}{16}$），
        只用八分音符（$\frac{2}{16}$）和
        十六分音符（$\frac{1}{16}$），且顺序有区别。例如
        八分音符--十六分音符--十六分音符，或
        十六分音符--八分音符--十六分音符。检验这两个例子是否合法，
        然后列出\emph{所有}可能的节奏。
        一共有多少种？
\end{enumerate}

\medskip
\noindent\textbf{第三部分——一个千年数列，以及
一个\omterm{def:g3:numbers:evenodd}{奇数}拍小节。}
\begin{enumerate}[resume]
  \item 下面来数一数，鼓手在更长音符上的节奏有多少种。
        每种节奏要么以十六分音符结尾，要么以
        八分音符结尾。请解释这为何给出如下规律：（填满
        $n$ 个十六分音符的节奏数）$=$（填满 $n - 1$ 个十六分音符的节奏数）
        $+$（填满 $n - 2$ 个十六分音符的节奏数）。
        以填满一个十六分音符的 $1$ 种节奏、填满两个的 $2$ 种节奏为起点，
        列出表格直至 $n = 8$：填满一个二分音符的节奏有多少种？
  \item 你算出的数——$1, 2, 3, 5, 8, 13, 21, 34$
——约在 1150 年由印度学者赫马昌德拉发表，
        他正是在计数这类节奏，
        比斐波那契在意大利将其写下早了两代。
        用一句话说明这些数的定义规律，
        并在表中最后三个数上加以验证。
  \item “摇摆”风格将两个等长的八分音符替换为
        一组长--短音对：先为 $\frac16$，随后为 $\frac{1}{12}$。
        请验证这一替换是否合法（总时值相同），并利用
        公\omterm{def:g4:fractions:def}{分母} $12$ 进行计算。
  \item \emph{三连音}将三个等长的音符压缩进一个
        四分音符中。每个音符的时值是多少（可将答案乘以 $3$ 来检验，
        见\cref{prop:g7:fractions:mult}）？
        对于三个等长音符填满一个二分音符的情形，同样的问题
        又该如何回答？
  \item 巴尔干舞曲采用七八拍：每小节
        的总时值为 $\frac78$。验证附点四分音符 $+$ 四分音符
        $+$ 四分音符恰好能填满它。接着解释为什么
        仅用二分音符和四分音符的任何组合，
        无论用多少个，都\emph{绝不}能填满一个七八拍小节。
        （以八分音符为单位计数，并思考\omterm{def:g3:numbers:evenodd}{偶数}与\omterm{def:g3:numbers:evenodd}{奇数}的关系。）
\end{enumerate}
\end{problem}
